% ===========================================================================
\chapter{Safe rules}\label{ch:safety}

Not every rule makes sense. Consider

\begin{shell}
likes(X, Y) :- person(X).
\end{shell}

``Everyone likes everything.'' The head has a variable, \code{Y}, that the
body never gives a value, so the rule would derive \code{likes(ann, Y)}: a
``fact'' with a variable in it, which is not a fact at all. The argument in
Chapter~\ref{ch:fixpoint} that evaluation stops also breaks: it relied on
derived facts using only constants that already exist.

A test has a similar problem. In \code{p(X) :- X != a, q(X).} the test comes
first, when \code{X} has no value yet, and there is nothing to compare.

Rules that avoid both problems are called \emph{safe}, and Datalog only
accepts safe rules. Our definition is:
\begin{itemize}
  \item every variable in the head appears in a positive atom of the body;
  \item every variable in a test or a negated atom appears in a positive
    atom \emph{before} it.
\end{itemize}
The second is a little stricter than it needs to be, since an engine could
reorder the body itself. Asking the programmer to write tests after the
atoms that bind their variables keeps the join of Chapter~\ref{ch:joins}
simple, and it is how most people write rules anyway.

\section{Checking a rule}

To check the second condition we need the variables of a literal, which for
a test means the variables among its two terms:

\begin{addbefore}{terms.nex}{before the Rule class}{-- head :- body. A rule with an empty body would be a fact.}
-- The variables a literal mentions.
function literal_variables(l: Literal): Array[String] do
  match l of
    Positive(atom) then result := atom.variables
    Compare(op, left, right) then
      do
        result := []
        across [left, right] as t do
          let name := variable_name(t)
          if name /= "" and not result.contains(name) then
            result.add(name)
          end
        end
      end
    Negated(atom) then result := atom.variables
  end
end
\end{addbefore}

A rule checks itself by walking its body in order, collecting the variables
that positive atoms bind, and complaining about the first variable used
before it is bound. It returns the complaint, or \code{""}, so that an error
message can say exactly what is wrong.

\begin{addbefore}{terms.nex}{to Rule, before text}{  text: String do|    result := head.text + " :-"}
  -- Why the rule is not safe, or "" if it is. A rule is safe when every
  -- variable is bound by a positive literal: those in the head by one
  -- anywhere in the body, those in a test by one before it.
  safety_problem: String do
    result := ""
    let bound: Array[String] := []
    across body as l do
      match l of
        Positive(atom) then
          across atom.variables as v do
            if not bound.contains(v) then
              bound.add(v)
            end
          end
        _ then
          across literal_variables(l) as v do
            if result = "" and not bound.contains(v) then
              result := "variable " + v + " in '" + literal_text(l)
                        + "' must appear in a positive literal before it"
            end
          end
      end
    end
    across head.variables as v do
      if result = "" and not bound.contains(v) then
        result := "variable " + v + " in the head must appear in a positive literal of the body"
      end
    end
  end

  is_safe: Boolean do
    result := safety_problem = ""
  end
\end{addbefore}

\begin{append}{terms.nex}
-- Are all of `rules` safe?
function all_safe(rules: Array[Rule]): Boolean do
  result := true
  across rules as r do
    if not r.is_safe then
      result := false
    end
  end
end
\end{append}

\section{A precondition}

Evaluating an unsafe rule is a mistake by whoever called the evaluator, so
it is a precondition of \code{run}, which every evaluator inherits:

\begin{tightaddbefore}{evaluate.nex}{to run, before deferred}{  deferred|  ensure|    fixpoint: is_fixpoint(rules, db)}
  require
    safe: all_safe(rules)
\end{tightaddbefore}

The engine's own command line will check every rule before it evaluates
anything, and report an unsafe one with its message. The precondition is
there for code that forgets.

\section{Trying it out}

\begin{tryit}{try_safety.nex}
intern Parser
intern Evaluate

match parse_program("p(X) :- q(X), r(X, Y).
p(X) :- q(Y).
p(X) :- X != a, q(X).
p(X) :- q(X), not r(Y).") of
  Ok as ok then
    do
      across ok.value.rules as r do
        if r.is_safe then
          console.print_line("safe:    " + r.text)
        else
          console.print_line("unsafe:  " + r.text)
          console.print_line("         " + r.safety_problem)
        end
      end
      do
        let e := create Naive_Evaluator.make
        e.run(ok.value.rules, create Database.make)
      rescue
        console.print_line(exception.to_string)
      end
    end
  Err as e then console.print_line(e.error)
end
\end{tryit}
\tryoutput{try_safety}
